\scanpage{303}
\mspage{44}
la \uline{portée} de $f$ sur $u$. Alors \deleted{$\unsure{\tau|E}$} $(\tau',f)$ est déterminé
de façon unique par le couple
\[
 (E,\pi),\quad E\in\unsure{\mathscr{P}|\mathcal{P}}(\mathrm{Tr}(\tau)),\quad
 \underset{\text{\(\hookrightarrow\) pondération}}{\pi}\colon E\longrightarrow
 \underset{\substack{\parallel\\\text{ens. des}\\\text{entiers }\geq 1}}{\mathbb N^*}.
\]
Ou si on préfère, par l'application $\pi=\pi_f$
\[
 \mathrm{Tr}(\tau)\xrightarrow{\pi}\mathbb N,\quad
 u\longmapsto\pi(u)=\pi(u,f)=(\inserted{\operatorname{long}}\pi_{\tau'}(f(u)))
\]
(qui redonne $E$ comme
\[
 E=\pi^{-1}(\mathbb N^*)\;).
\]
Cette application n'est soumise à aucune restriction
\[
 \operatorname{Ob}\,{}_{\tau\backslash}\underline{\mathrm{Ty}}\deleted{\unsure{/\tau}}
 \xrightarrow[\ ]{\overset{\mathrm{bij.}}{\sim}}
 \deleted{\unsure{\mathrm{Appl.}|\mathrm{Appl}}}\operatorname{Hom}_{\mathrm{Ens}}(\mathrm{Tr}(\tau),\mathbb N)
\]
On récupère $\tau'$ à partir de $\pi$ ainsi. On
prend $E=\pi^{-1}(\mathbb N^*)$ (fl. de transition \deleted{\unsure{pour}} \inserted{non dégénérées}
pour $\pi$), $E$ est \deleted{\unsure{totalemt|totalement}} ordonné par l'ordre induit
par celui de $\mathrm{Tr}(\tau)$, il est de plus signaturé
par les signatures de $\mathrm{Tr}(\tau)$, soient
\[
 (\varepsilon_1,n_1)\mathbin{-}(\varepsilon_r,n_r)
\]
la suite des signatures et pondérations, \deleted{associées} \inserted{de la}
suite des éléments de $E$. \quad ($\varepsilon_i\in\{\pm1\}$, $n_i\in\mathbb N^*$
pour $i\in[1,r]$). Alors,
\[
 \tau'=\Delta_{\varepsilon_1}^{\deleted{\unsure{\varepsilon_1}}n_1}\Delta_{\varepsilon_2}^{n_2}
 \deleted{\unsure{\Delta_{\varepsilon_3}^{n_3}}}\ldots\Delta_{\varepsilon_r}^{n_r}
 \qquad r=\operatorname{card}E
\]
dans le monoïde libre $\mathrm{Ty}$, où on a désigné
par $\Delta_+,\Delta_-$ (au lieu de $\Delta^+,\Delta^-$) les
deux générateurs. \unsure{et application|et l'application}
\[
 I_\tau\unsure{\mathbin{-}|\longrightarrow} I_{\tau'}
\]
% archive sheet number: 302
